\section{Rapid decay in the adelic quotient height}

% BEGIN R89D1
\paragraph{R89D1. The quotient height on Siegel representatives.}
Fix the Hermitian space, the prescribed infinite embeddings and a
compact open finite level $K$ for $G=U(V)$. Write $H$ for the source's
matrix height and $H_{\rm quot}(g)=\inf_{\gamma\in G(F)}H(\gamma g)$.
In each of the finitely many finite components of FV9--10 choose $s$.
Choose the real arithmetic group $\Gamma_s$ of RPI1, the Siegel set
$\mathcal S$ and its finite set of rational translates $q$ of RPI8.
Every $g$ has a left rational translate $h$ of the form
\[
               h_f\in sK,\qquad h_\infty=qz,\quad z\in\mathcal S.
\]
With $f$ the number of infinite factors and the faithful matrix norm
$\|\cdot\|_{\rm lin}\ge1$ of RPI1, on this representative one has
\[
             H_{\rm quot}(g)\le H(h)\le C\|z\|_{\rm lin}^{f}.
                                                               \tag{D89.1}
\]
The constant is independent of $g$ on all these components.

Indeed first put the finite component in $sK$ by left rational
translation. Reduce the real component further by $\Gamma_s$.
Each element used in this second reduction is the infinite component
of a rational element whose finite component preserves $sK$ on the
left. The finite-component condition thus survives. The finite union
of rational translates of $\mathcal S$ is the usual reduction theorem,
in Miller--Schmid, \emph{On the rapid decay of cuspidal automorphic
forms}, (2.8), applied in RPI8 to $\operatorname{Res}_{F/\mathbb Q}G$.

The product defining $H$ measures both a matrix and its inverse. On
$sK$ the finite matrix entries of both, expressed in the specified
basis of the coefficient extension, vary in compact subsets of the
finite adeles. Each such compact set has integral coordinates at all
but a fixed finite set of places. Cover it by basic restricted-product
neighbourhoods and take a finite subcover. At the remaining places
the coordinates are bounded. Thus the finite product is bounded on
$sK$. At each infinite factor the maximum of $1$, the matrix
operator norm, and the inverse operator norm is at most a constant
times $\|qz\|_{\rm lin}$, by the standard and dual matrix
summands used in RPI1. This norm is submultiplicative and there are
only finitely many $s,q$. Multiplication over the $f$ factors proves
(D89.1), including the compact factors of the Siegel set in the norm
of the whole element $z$. The first inequality is the definition of
the infimum and uses no minimum or converse height comparison.
% END R89D1

% BEGIN R89D2
\paragraph{R89D2. The original rapid-decay input, with every quotient-height exponent.}
Let $\phi$ satisfy the original automorphy, smoothness, uniform
moderate growth and cuspidality hypotheses for this $D,V,\theta$ and
its given exponential family. For every integer $n\ge0$ there is a
constant $B_{\phi,n}$ such that
\[
        |\phi(g)|\, H_{\rm quot}(g)^n\le B_{\phi,n}
                        \qquad(g\in G(\mathbb A_F)).
\]
This is the full predicate $\mathrm{Published.CuspRapidDecay}$ used in
the construction of the spectral comparison.

Choose $K$ in the given finite right level. On a finite component put
$f_s(u)=\phi(su)$. The exponential family gives the actual smooth
right derivatives on $G_\infty$ as in RPI1. Uniform moderate growth
already bounds every differential word of $\phi$ by $C_D H^{N_0}$
with one exponent $N_0$. At $su$ the finite height factor is fixed
and the infinite product is at most $C\|u\|_{\rm lin}^{f}$ by the
same operator-norm comparison as in R89D1. Hence all the derivatives
of $f_s$ have the single growth exponent $fN_0$. Rational left
invariance and right $K$-invariance make $f_s$ invariant under
$\Gamma_s$. RPI2--6 apply to these smooth components. Splitting the
unipotent quotient at the finite level gives the vanishing of the
real constant terms for the prescribed isotropic subspaces, and
integration in stages gives vanishing for every proper rational
parabolic. The same vanishing holds for all right derivatives.

For the real decay one may now cite Borel--Labesse--Schwermer,
\emph{On the cuspidal cohomology of $S$-arithmetic subgroups of
reductive groups over number fields}, \S\S2.1, 4.4--4.5.
Their Arthur truncation $\Lambda^T f$ of a smooth function of uniform
moderate growth is rapidly decreasing on the norm-one group. The
difference $f-\Lambda^T f$ is a sum indexed by the proper parabolics,
each term built from a constant term of $f$. All these terms vanish
here. We apply their arithmetic case to $\operatorname{Res}_{F/\mathbb Q}G$
with $S$ consisting of the infinite place. The real centre is
compact, so the norm-one group is the full group. Apply the result
also to $u\mapsto f_s(qu)$ on $q^{-1}\Gamma_s q$. Rational conjugation
permutes the rational parabolics and changes their Haar measures by
constants, and the same derivative exponent works. Consequently
$|f_s(qz)|\le C_{s,q,M}\|z\|_{\rm lin}^{-M}$ on $\mathcal S$
for every $M\ge0$. Equivalently, (RPI.28) proves this estimate
directly by root contraction in the same norm, including the compact
unipotent and maximal compact factors. The cited formulation applies
without finiteness under the maximal compact group or the
enveloping-algebra centre. In rational rank zero the Siegel set is
compact and continuity suffices.

For a representative $h$ in R89D1, automorphy and finite-level
invariance give $\phi(g)=\phi(h)=f_s(qz)$. Set $M=fn$ and use
(D89.1). The product on the left of the claimed bound is at most
$C^n\max_{s,q} C_{s,q,fn}$, which is independent of $g$. Rank zero
for $V$ itself is immediate. This proves the statement at arbitrary
$D$ with the required order of the quantifiers. The five inputs in the
linked-bundle construction include this decay theorem. The corresponding
finiteness input has a separate classical proof below.
% END R89D2

% BEGIN R89D3
\paragraph{R89D3. The fixed-type finiteness theorem in the precise coordinates used here.}
Fix the datum, embeddings and Hermitian space, the source's open
level subgroup, a frame and finite-dimensional continuous compact
representation at every infinite place, and the specified list
$c_{\tau,k}$ of eigenvalues. The span of all the smooth automorphic
functions satisfying those finite-level, compact-type and differential
equations is finite dimensional. This proves the full input
$\mathrm{Published.HarishChandraFinite}$, including functions with
nonzero constant terms.

Take a compact open subgroup $K$ of the finite slice of the given
level. Restriction to the finitely many components $s$ of FV9--10
embeds the space of functions into the direct sum of spaces on
$\Gamma_s\backslash G_\infty$. Each component function is smooth
and has moderate growth by the empty differential word of its
growth condition and the matrix-norm comparison of RPI1. The
frames identify the compact groups in the hypothesis with maximal
compact subgroups at the real places. Let $\xi_\tau$ be the central
compact convolution idempotent projecting onto the finite list of
irreducible types in the given representation at $\tau$. On every
right compact orbit the function is a linear combination of matrix
coefficients of this representation, by the prescribed type.
Convolution with $\xi_\tau$ fixes every such coefficient. Evaluating
at the identity of the compact orbit shows the function itself is fixed.
The separate-place projections commute, and their product is one
fixed finite-type idempotent $\xi$ for all the component functions.

It remains to check which part of the centre the displayed differential
equations fix. Put $r=\dim_E V>0$. At each real place the complexified
Lie algebra is $\mathfrak{gl}_r$, in the given matrix coordinates.
The definition of the complex derivative decomposes a matrix into its
two real Lie-algebra parts using $Z\mapsto-H^{-1} Z^* H$. Thus the
operator indexed by $k$ is the action of the cyclic product
\[
 T_{k+1}=\sum_{i_1,\ldots,i_{k+1}=1}^r
    E_{i_1 i_2} E_{i_2 i_3}\cdots E_{i_{k+1}i_1}
       \quad(T_1=\textstyle\sum_i E_{ii})
\]
in the universal enveloping algebra. Commuting any matrix unit
through a cyclic product, the two terms at each contracted index
cancel after summation. Hence all $T_j$ are central. Their leading
symbols in the degree filtration are $\operatorname{tr}(X^j)$.
For $1\le j\le r$ these generate the polynomial invariants for
conjugation: restriction to a diagonal matrix gives the power sums,
which generate the symmetric polynomials by Newton's identities,
and matrices with distinct eigenvalues are dense. The leading
symbol of any central element is invariant, by the commutator
formula for the degree filtration. Subtracting a polynomial of the
same leading degree in $T_1,\ldots,T_r$ and inducting on the degree
proves that they generate the whole centre.

It follows that the equations $T_j=c_{\tau,j-1}$, at every place
and for $1\le j\le r$, put all the component functions in the
kernel of a single ideal of finite codimension in the centre of the
product Lie algebra. The quotient is generated by scalars and
has dimension at most one. The additional equations can only
restrict the space. Harish-Chandra's finiteness theorem in
A.~Borel, \emph{Automorphic Forms on Reductive Groups},
Theorem 8.4, gives finite dimension on each arithmetic
quotient for this ideal and $\xi$. The finite direct sum gives the
assertion. For $r=0$ the group is trivial and the functions have
dimension at most one. This theorem and R89D2 supply two of the
five real-group inputs to the linked-bundle period comparison.
% END R89D3
